\documentclass[10pt]{article}

%====================================
%				package
%====================================

%gestion du français et des accents
\usepackage[utf8]{inputenc}
\usepackage[T1]{fontenc}
\usepackage[french]{babel}

%symbole mathématique
\usepackage{amsmath}
\usepackage{amsfonts}
\usepackage{bbm}
\usepackage{amsthm}
\usepackage{amssymb}
\usepackage[colorlinks=true, linkcolor=blue, urlcolor=blue]{hyperref}
\usepackage{bm} 
\usepackage{stmaryrd}
\usepackage{dsfont}
\usepackage{mathtools} % pour mathrlap


%mise en page
\usepackage[margin=1in]{geometry}
\usepackage{lmodern}
\usepackage{fancyhdr}
\usepackage{soul}
\usepackage{float}
\usepackage{titlesec}
\usepackage{enumitem}
\usepackage{tcolorbox} 
\usepackage{tabularx}
\usepackage{stmaryrd}
\usepackage{pgfplots}
\pgfplotsset{compat=1.18}
\usepackage{tikz}
\usepackage{xcolor}
\pagestyle{fancy}    % Utiliser un style d'en-têtes personnalisés
\fancyhf{}           % Réinitialise les en-têtes et pieds de page
\fancyhead[L]{Anthony Gabino}  % Nom de l'auteur à gauche
\fancyhead[R]{Exercice à rendre semaine 11}  % Titre du sous-chapitre à droite
\fancyfoot[C]{\thepage}  % Numéro de page centré en bas
\begin{document}
	

\section*{Exercice à rendre}
Soit $A$ l'ensemble définit par
\[
A = \left \{ (x,y) \in ]0,1] \times \mathbb R_+ ~|~ 0 \leqslant y \leqslant 2 + \sin \frac1x \right\}
\]
On remarque alors que $A$ décrit la portion d'aire sous le graphe de la fonction $x \mapsto 2 + \sin 1/x$ pour $x \in ]0,1]$, ainsi on voit que
\[
\partial A = {\color{orange} \mathcal G (2 + \sin 1/x)} \cup {\color{red}\left([0,1] \times \{0\} \right)}\cup {\color{blue}\left( \{1\}\times[0,2 +\sin1]\right)} \cup {\color{green!50!black}(\{0\}\times [0,3])}
\]
\begin{figure}[H]
\begin{center}
\begin{tikzpicture}
	\begin{axis}[
		width=8cm,
		height=5cm,
		xmin=0,
		xmax=1,
		ymin=0,
		ymax=3.5,
		axis lines=middle,
		xlabel={$x$},
		ylabel={$y$},
		samples=1200,
		domain=0.001:1,
		clip=false
		]
		
		% Aire correspondant à A
		\addplot[
		draw=none,
		fill=blue!25,
		]
		{2 + sin(deg(1/x))}
		\closedcycle;
		
		% Bord supérieur plus foncé
		\addplot[
		very thick,
		orange,
		]
		{2 + sin(deg(1/x))};
		
		% Bord inférieur y=0
		\addplot[
		very thick,
		red,
		domain=0:1
		]
		{0};
		
		% Bord vertical en x=0
		\draw[very thick, green!50!black]
		(axis cs:0,0) -- (axis cs:0,3);
		
		% Bord vertical en x=1
		\pgfmathsetmacro{\h}{2 + sin(deg(1))}
		
		\draw[
		very thick,
		blue!80!black
		]
		(axis cs:1,0) -- (axis cs:1,\h);
		
		
	\end{axis}
\end{tikzpicture}
\caption{Graphe de la fonction $x \mapsto 2 + \sin \frac1x$.}
\end{center}
\end{figure}
Or il est claire que les segments sont tous négligeable, il reste donc à montrer que le graphe ${\color{orange} \mathcal G}$ est aussi de négligeable. Premièrement, on remarque que la fonction $f : x \mapsto 2 + \sin \frac1x$ est intégrable au sens de Riemann sur $]0,1]$, puisque continue sur cette intervalle, ainsi pour tout $\varepsilon > 0$, il existe une partition $P$ de $[0,1]$ telle que, si on note $\tilde f$ le prolongement par zéro de $f$, on a
\[
\sum_{Q \in P} \sup_{x \in Q} \tilde f(x) \mathrm{Vol}(Q) - \sum_{Q \in P} \inf_{x \in Q} \tilde f(x) \mathrm{Vol}(Q) \leqslant \varepsilon
\]
D'où
\[
\sum_{Q \in P} \mathrm{Vol} \left(Q \times \left[\inf_{x \in Q} \tilde f(x), \sup_{x \in Q} \tilde f(x) \right]\right) \leqslant \varepsilon
\]
Or par définition
\[
\mathcal G(f) \subset \mathcal G(\tilde f) \subset \bigcup_{Q \in P} \left(Q \times \left[\inf_{x \in Q} \tilde f(x), \sup_{x \in Q} \tilde f(x) \right]\right) \leqslant \varepsilon
\]
Ainsi, pour tout $\varepsilon > 0$, on peut trouver une famille finie de pavés dont l'union contient $\mathcal G(f)$ et dont la somme des volumes et strictement inférieure à $\varepsilon$, ainsi $\mathcal G(f)$ est bien négligeable. 

Or comme la frontière $\partial A$ est l'union finie d'ensemble négligeable, alors elle l'est aussi. Donc on conclut que $A$ est bien mesurable au sens de Jordan. 
\end{document}